\section{연습문제}
\label{sec:finite-field-exercises}

문제는 A, B, C의 세 단계로 나뉜다. 별표 문제는 다음 두 절에 상세 풀이가 있다.

\subsection*{A. 판정과 계산}

\begin{exercise}
원소가 다음 개수인 체의 존재 여부를 판정하라.
\[
  6,\quad8,\quad9,\quad12,\quad25,\quad27.
\]
존재하는 경우 표수와 소체 위의 차수도 구하라.
\end{exercise}

\begin{exercise}[*]
$\alpha^2+\alpha+1=0$인 $\alpha$를 이용하여 $\F_4=\F_2(\alpha)$를 구성하라.
\begin{enumerate}[label=(\alph*)]
  \item 네 원소를 모두 적어라.
  \item 덧셈표와 곱셈표를 작성하라.
  \item 각 원소가 $X^4-X$의 근임을 확인하라.
\end{enumerate}
\end{exercise}

\begin{exercise}[*]
$\alpha^3+\alpha+1=0$인 $\alpha$에 대해 $\F_8=\F_2(\alpha)$라 하자.
\begin{enumerate}[label=(\alph*)]
  \item $1,\alpha,\dots,\alpha^6$을 $1,\alpha,\alpha^2$의 선형결합으로 나타내라.
  \item $\alpha$가 $\F_8^\times$의 생성원임을 보여라.
  \item $(\alpha^2+\alpha+1)^{-1}$을 구하라.
\end{enumerate}
\end{exercise}

\begin{exercise}
$\F_9=\F_3(i)$, $i^2=-1$이라 하자.
\begin{enumerate}[label=(\alph*)]
  \item $(1+i)^k$를 $1\le k\le8$에 대해 계산하라.
  \item $1+i$가 원시원임을 보여라.
  \item $(2+i)^{-1}$을 구하라.
\end{enumerate}
\end{exercise}

\begin{exercise}[*]
$\F_2[X]$에서 $X^4-X$를 일계수 기약다항식의 곱으로 분해하라. 이 분해가 $\F_4$의 원소들과 어떻게 연결되는지 설명하라.
\end{exercise}

\begin{exercise}
다음 포함관계가 성립하는지 판정하라.
\[
\F_{2^2}\subseteq\F_{2^6},\quad
\F_{2^3}\subseteq\F_{2^{12}},\quad
\F_{3^4}\subseteq\F_{3^{10}},\quad
\F_{5^6}\subseteq\F_{5^{18}}.
\]
포함이 성립하면 확장차수도 구하라.
\end{exercise}

\begin{exercise}[*]
고정된 $\overline{\F_p}$ 안에서
\[
  \F_{p^6}\cap\F_{p^{10}}
  \quad\text{및}\quad
  \F_{p^6}\F_{p^{10}}
\]
을 구하라. 두 체의 합성체 위에서 각 체의 차수도 구하라.
\end{exercise}

\begin{exercise}
$L=\F_{3^4}$라 하자.
\begin{enumerate}[label=(\alph*)]
  \item $\Aut_{\F_3}(L)$의 원소를 모두 상대 프로베니우스의 거듭제곱으로 적어라.
  \item 각 비자명한 자동동형의 위수를 구하라.
  \item 각 자동동형의 고정체를 구하라.
\end{enumerate}
\end{exercise}

\begin{exercise}[*]
$\F_8=\F_2(\alpha)$, $\alpha^3+\alpha+1=0$이라 하자.
\begin{enumerate}[label=(\alph*)]
  \item $\alpha$의 프로베니우스 궤도를 구하라.
  \item 궤도의 곱으로 $\alpha$의 최소다항식을 복원하라.
  \item $\alpha+1$의 프로베니우스 궤도와 최소다항식을 구하라.
\end{enumerate}
\end{exercise}

\begin{exercise}[*]
공식
\[
  N_q(n)=\frac1n\sum_{d\mid n}\mu(d)q^{n/d}
\]
을 이용하여
\[
  N_2(1),\ N_2(2),\ N_2(3),\ N_2(4),\ N_3(2)
\]
를 계산하라.
\end{exercise}

\subsection*{B. 핵심 정리의 증명}

\begin{exercise}
유한체 $F$의 원소 수가 반드시 $p^n$ 꼴임을 표수와 벡터공간 차원을 사용하여 증명하라.
\end{exercise}

\begin{exercise}[*]
$q=p^n$이라 하고 $\overline{\F_p}$ 안에서
\[
  S=\{a:a^q=a\}
\]
라 하자. $S$가 덧셈, 곱셈과 역원에 닫혀 있는 체이며 정확히 $q$개의 원소를 가짐을 증명하라.
\end{exercise}

\begin{exercise}
원소가 $q=p^n$개인 모든 체가 $X^q-X$의 $\F_p$ 위 분해체임을 보여 유한체의 유일성을 증명하라.
\end{exercise}

\begin{exercise}[*]
고정된 $\overline{\F_p}$ 안에서
\[
  \F_{p^m}\subseteq\F_{p^n}
  \quad\Longleftrightarrow\quad
  m\mid n
\]
을 증명하라.
\end{exercise}

\begin{exercise}
앞 문제만을 이용하여
\[
  \F_{p^m}\cap\F_{p^n}=\F_{p^{\gcd(m,n)}}
\]
과
\[
  \F_{p^m}\F_{p^n}=\F_{p^{\operatorname{lcm}(m,n)}}
\]
을 증명하라.
\end{exercise}

\begin{exercise}
유한체의 모든 대수적 확장이 정규이고 분리임을 증명하라. 유한확장과 무한 대수적 확장의 경우를 구분하여 논증하라.
\end{exercise}

\begin{exercise}
체의 곱셈군의 유한부분군이 순환군이라는 정리를 사용하여 $\F_q^\times$가 위수 $q-1$인 순환군임을 증명하라. 이어서 원시원이 $\varphi(q-1)$개임을 보여라.
\end{exercise}

\begin{exercise}[*]
$L=\F_{q^n}$라 하자. 상대 프로베니우스
\[
  \Frob_q:L\to L,\qquad a\mapsto a^q
\]
가 $\F_q$-자동동형이며 위수가 정확히 $n$임을 증명하라. 이어서
\[
  \Aut_{\F_q}(L)=\langle\Frob_q\rangle
\]
임을 보여라.
\end{exercise}

\begin{exercise}
$L=\F_{q^n}$에서 $\Frob_q^d$의 고정체가
\[
  \F_{q^{\gcd(n,d)}}
\]
임을 증명하라.
\end{exercise}

\begin{exercise}[*]
$f\in\F_q[X]$가 일계수 기약 $d$차다항식이라 하자. 다음 동치를 증명하라.
\[
  f\mid X^{q^n}-X
  \quad\Longleftrightarrow\quad
  d\mid n.
\]
\end{exercise}

\subsection*{C. 구조와 응용}

\begin{exercise}
$\F_q[X]$에서
\[
  X^{q^n}-X
  =\prod_{\substack{f\text{ monic irreducible}\\\deg f\mid n}}f(X)
\]
임을 증명하고 각 인수가 한 번만 나타나는 이유를 설명하라.
\end{exercise}

\begin{exercise}[*]
$N_q(n)$을 $\F_q[X]$의 일계수 기약 $n$차다항식의 개수라 하자.
\begin{enumerate}[label=(\alph*)]
  \item $q^n=\sum_{d\mid n}dN_q(d)$를 증명하라.
  \item 모비우스 반전으로
  \[
    N_q(n)=\frac1n\sum_{d\mid n}\mu(d)q^{n/d}
  \]
  를 유도하라.
\end{enumerate}
\end{exercise}

\begin{exercise}[*]
고정된 $\overline{\F_p}$ 안에서
\[
  K=\bigcup_{n\ge1}\F_{p^{n!}}
\]
라 하자. $K$가 체이고 $\F_p$ 위에서 대수적이며 대수적으로 닫혀 있음을 증명하라.
\end{exercise}

\begin{exercise}
$\F_{p^n}$의 부분체가 $n$의 양의 약수와 일대일로 대응함을 증명하라. 부분체 포함관계가 약수의 나눗셈 관계와 같은 방향임을 설명하라.
\end{exercise}

\begin{exercise}[*]
$L=\F_{q^{12}}$와 $G=\Aut_{\F_q}(L)=\langle\Frob_q\rangle$를 생각하자.
\begin{enumerate}[label=(\alph*)]
  \item $G$의 부분군을 모두 적어라.
  \item 각 부분군의 고정체를 구하라.
  \item 부분군 포함과 고정체 포함의 방향이 반대임을 확인하라.
\end{enumerate}
\end{exercise}

\begin{exercise}
$f=X^4+X+1\in\F_2[X]$라 하자.
\begin{enumerate}[label=(\alph*)]
  \item $f$가 기약임을 보여라.
  \item $\F_{16}=\F_2[X]/(f)$를 구성하라.
  \item $\F_{16}$의 모든 부분체와 $\F_2$-자동동형을 구하라.
\end{enumerate}
\end{exercise}

\begin{exercise}
한 체 안의 두 유한부분체가 같은 수의 원소를 가지면 서로 같음을 $X^q-X$의 근 집합을 이용하여 증명하라.
\end{exercise}

\begin{exercise}[*]
유한체 $F$의 모든 영이 아닌 원소의 곱이 $-1$임을 역원끼리 짝지어 증명하라. 이를 $F=\F_p$에 적용하여 윌슨 정리를 유도하라.
\end{exercise}

\begin{exercise}
$\F_2[X]$의 일계수 기약 육차다항식의 개수 $N_2(6)$을 구하라. 이어서 $X^{64}-X$에 나타나는 기약인수들의 차수별 개수를 정리하라.
\end{exercise}

\begin{exercise}[*]
$\alpha\in\F_{q^n}$이고 $d$가 $\alpha^{q^d}=\alpha$를 만족하는 최소의 양의 정수라 하자.
\begin{enumerate}[label=(\alph*)]
  \item $d\mid n$임을 보여라.
  \item
  \[
    \prod_{j=0}^{d-1}(X-\alpha^{q^j})
  \]
  의 계수가 $\F_q$에 속함을 증명하라.
  \item 이 곱이 $\alpha$의 $\F_q$ 위 최소다항식임을 증명하라.
\end{enumerate}
\end{exercise}
